\section{Common scalar recovery of two marked coefficients}

These are the interpolation and ordered-word identities used in
\cite[Lemma~4.2]{Arad2013AreaLaw} and
\cite[Proposition~5.7]{Arad2016RigorousRG}.
Only a spanning family of evaluations is required; the coefficient operators
need not be linearly independent. The distinction between spanning and a
basis in the source interpolation argument is recorded in
\cite{gap:aklv13_interpolation_and_word_count}.

\begin{definition}[Ordinary labels]
    \label{def:two_marker_ordinary_labels}
    \lean{MvPolynomial.twoMarkerOrdinaryLabels}
    For a finite alphabet $I$ and distinguished labels $u,v\in I$, let
    \begin{align}
        \mathcal O_{u,v}=I\setminus\{u,v\}.
        \label{eq:two_marker_ordinary_labels}
    \end{align}
    The definition does not require the two labels to be distinct.
\end{definition}

\begin{definition}[The full-frequency ordered-word coefficient]
    \label{def:full_frequency_word_coefficient}
    \lean{WordMultiplicity.wordCoefficient}
    Let $I$ be finite, let $R$ be a semiring, and let $H_a\in R$ for $a\in I$.
    For $r\ge0$ and a count vector $\nu:I\to\N$, put
    \begin{align}
        \mathcal W_{H,r}(\nu)=
        \sum_{\substack{w:\{0,\ldots,r-1\}\to I \\
                |w^{-1}(a)|=\nu(a)\ (a\in I)}}
        H_{w(0)}H_{w(1)}\cdots H_{w(r-1)}.
        \label{eq:full_frequency_word_coefficient}
    \end{align}
    The factors retain their displayed order. An empty product is $1$, and a
    sum with no words is zero.
\end{definition}

\begin{theorem}[One sample family for every degree and frequency slice]
    \label{thm:two_window_common_scalar_recovery}
    \lean{Polynomial.exists_twoWindow_wordCoefficient_recovery}
    \uses{def:two_marker_ordinary_labels,def:full_frequency_word_coefficient}
    Let $\ell,K$ be nonnegative integers, let
    \begin{align}
        I=(\{0,\ldots,2\ell-1\}\sqcup\{0,\ldots,2\ell-1\})
        \sqcup\{0,1,2\},
        \label{eq:two_window_scalar_alphabet}
    \end{align}
    and choose $0\le p,s<2\ell$.
    Let $u$ be the label $p$ in the first window and $v$ the label $s$ in the
    second window. Let $R$ be a ring and a complex algebra, without any
    commutativity assumption on $R$, and put $N=(K+1)^{4\ell+1}$.

    There are complex sample points $z_t:I\to\C$, $0\le t<N$, with the
    following property. For every $i,j\ge0$ and every scalar mask
    $\mu:(I\to\N)\to\C$, there are complex coefficients $\gamma_t$ such
    that, for every family $H:I\to R$ and every $0\le r\le K$,
    \begin{align}
         & \sum_{\substack{\nu:I\to\N         \\ \sum_a\nu(a)=r\\
                \nu(u)=i,\ \nu(v)=j}}
        \mu(\nu)\,\mathcal W_{H,r}(\nu)\notag \\
         & \qquad=\sum_{t=0}^{N-1}\gamma_t
        [x^iy^j]\left(D_t+yH_v+xH_u\right)^r,
        \qquad D_t=\sum_{a\in\mathcal O_{u,v}}z_t(a)H_a.
        \label{eq:two_window_common_scalar_recovery}
    \end{align}
    Here $x,y$ are central formal variables and scalar products mean the
    complex algebra action. The selected terms $H_u,H_v$ remain unweighted.
    The points precede $i,j,\mu$, while the recovery coefficients precede $H$
    and every degree $r\le K$. Points for different selected pairs need not
    coincide. The case $K=0$ is included.
\end{theorem}
\begin{proof}
    \uses{def:two_marker_ordinary_labels,def:full_frequency_word_coefficient}
    The two distinguished labels are distinct, so the ordinary label set has
    $4\ell+1$ members. Ordinary count vectors with each coordinate between
    zero and $K$ form a monomial box of exactly $(K+1)^{4\ell+1}$ elements.
    Choose the finite interpolation points for this box before the frequencies
    or mask are supplied.

    Fix $i,j,\mu$. Restore the two distinguished counts in each ordinary box
    exponent, and use $\mu$ of that restored count vector as its scalar weight.
    Uniform scalar interpolation gives coefficients $\gamma_t$ before any
    module-valued coefficient family is supplied. A frequency slice at any
    $r\le K$ embeds in the same box; set the other coefficients to zero.
    Every count coordinate in that slice is at most $K$, so restoring and
    removing the distinguished counts are inverse on the retained slice.

    Apply this recovery identity to the coefficients
    $\mathcal W_{H,r}(\nu)$. The existing counted-word identity identifies
    each evaluated ordinary monomial sum with
    $[x^iy^j](D_t+yH_v+xH_u)^r$, without changing the order of the factors
    $H_a$. Reindexing the monomial box by its exact cardinality gives
    (\ref{eq:two_window_common_scalar_recovery}).
\end{proof}
